\chapter{Objetivos}
\label{ch:objetivos}
 

\paragraph*{}El objetivo principal de este Trabajo de Fin de Grado es ampliar y mejorar una aplicación interactiva orientada al análisis de supervivencia, aplicada al estudio del abandono académico en el ámbito universitario. No se trata solo de añadir nuevas funciones por añadirlas, sino de hacer que la herramienta sea más útil, más clara y más capaz de ayudar a interpretar un problema que, en la práctica, afecta tanto a los estudiantes como a las instituciones. Para ello, se plantea el desarrollo de nuevas funcionalidades que permitan analizar con mayor precisión los factores que influyen en este fenómeno.

\paragraph*{}La solución propuesta, al igual que en el trabajo previo, se desarrollará utilizando tecnologías como Python y Dash. A partir de estas herramientas, se busca construir una plataforma que facilite la realización de análisis interactivos, la visualización de resultados y la aplicación de técnicas estadísticas sobre datos educativos reales. Como punto de partida se utilizará el conjunto de datos empleado en el Trabajo de Fin de Grado anterior, lo que permitirá mantener una base común para comparar resultados y valorar con más justicia las mejoras introducidas.

\paragraph*{}Además, el proyecto no se limita únicamente al análisis del abandono académico. La idea es ir un paso más allá y mejorar la utilidad real de la herramienta desarrollada. En este sentido, se pretende evolucionar la aplicación hacia un sistema más completo, capaz de servir como apoyo en la toma de decisiones por parte de responsables académicos y de contribuir al diseño de estrategias orientadas a mejorar la retención del alumnado.

Para alcanzar este objetivo general, se plantean los siguientes objetivos específicos:

\begin{itemize}
    \item Implementar y probar una inteligencia artificial más rápida que la del trabajo previo, con el fin de optimizar los tiempos de generación y visualización de resultados. Para comprobar esta mejora, se medirán aspectos como el tiempo de respuesta, el tiempo de generación de resultados y el consumo de recursos.
    
    \item Incorporar soporte multilingüe, permitiendo que la aplicación pueda utilizarse en distintos idiomas, concretamente castellano e inglés. Además, se verificará que los elementos de la interfaz estén correctamente traducidos, de forma que la herramienta resulte más accesible para distintos perfiles de usuario.
    
    \item Ampliar el conjunto de variables utilizadas en el modelo, incorporando nuevos atributos y variables no binarias que permitan enriquecer el análisis y mejorar la capacidad explicativa de los modelos. Asimismo, se justificará la selección de dichas variables, explicando los criterios seguidos para incluirlas y valorando su relevancia frente a otros posibles atributos no considerados.
    
    \item Incluir gráficos y métodos estadísticos complementarios basados en el Test de Log-Rank y en el modelo de riesgos proporcionales de Cox. Con ello se pretende facilitar el análisis comparativo entre distintos grupos de estudiantes y ofrecer una interpretación más completa de los resultados.
    
    \item Explorar e integrar nuevas técnicas de análisis de supervivencia, como modelos paramétricos del tipo Weibull y Exponencial, junto con enfoques basados en aprendizaje automático como Random Survival Forest. La finalidad será comparar distintos enfoques, estudiar sus factores asociados y mejorar la precisión y versatilidad del sistema.
    
    \item Permitir la exportación automática de informes, imágenes y representaciones gráficas, facilitando la documentación de los resultados y su uso en contextos de investigación o gestión académica.
\end{itemize}

\paragraph*{}En conjunto, estos objetivos marcan el camino del proyecto y permiten comprobar, de forma ordenada, si las mejoras introducidas cumplen realmente con lo planteado en el anteproyecto. Además, este proceso de validación ayudará a evaluar la fiabilidad, la utilidad y el alcance de la herramienta desarrollada, evitando que la mejora quede solo en una ampliación técnica sin impacto práctico.
